\section{Proof}

The structure constants of $\cJ(m,k)$ are given by
\[
    p_{b,c}^a(m) = \sum_{i} \binom{k-a}{i}\binom{a}{k-b-i}\binom{a}{k-c-i} \binom{m-k-a}{b+c+i-k}
\]
(as in \cite{E,KK}).
Here $0 \leq a, b, c \leq k$, and $i$ can be restricted to the range
\[
    \max(0,k-a-b,k-b-c,k-a-c) \leq i \leq \min(k-a,k-b,k-c, m-a-b-c).
\]
In the following we always assume $m \geq 3k$. Under this assumption we have $p_{b,c}^a(m) > 0$ if and only if $a,b,c$ satisfy the triangle inequalities \cite[Lemma~4.1]{E}.
More precisely we have the following.

\begin{lemma}
    We have $p^a_{b, c}(m) > 0$ if and only if $|b - c| \leq a \leq b + c$. 
    Assuming that~$0 \leq b - c \leq a \leq b + c$, the leading term of $p^a_{b,c}(m)$ is
    \[
        \lt(p_{b,c}^a(m)) = \begin{cases}
        \binom{a}{b}\binom{a}{c}\frac{1}{(b+c - a)!} m^{b+c-a} &: a \geq b,\\
        \binom{k-a}{k-b}\binom{a}{b-c}\frac{1}{c!} m^{c} &: a\leq b.
        \end{cases}
    \]
In particular, $\deg(p_{b,c}^a(m))\leq \min(b,c)$, and equality holds if and only if $a \leq \max(b,c)$.
\end{lemma}

Now consider $\cJ(m, k)^d$.
For $a, b, c \in [0,k]^d$, let
\[
    p^a_{b,c}(m) = \prod_{i=1}^d p^{a_i}_{b_i,c_i}(m).
\]
These are the structure constants of $\cJ(m, k)^d$.
The following lemma generalizes the previous one.

\begin{lemma}
    Let $a,b,c\in[0,k]^d$.
    Then $p^a_{b,c}(m) > 0$ if and only if $|b-c| \leq a \leq b + c$
    Moreover $\deg(p_{b,c}^a(m)) \leq \wt(\min(b, c))$, with equality if and only if $a \leq \max(b, c)$.
    If~$\deg(p_{b,c}^a(m)) = \wt(b) = \wt(c)$ then $a \leq b = c$ and the leading term of $p^a_{b,c}(m)$ is
    \[
        \lambda(p^a_{b,c}(m))
        = \binom{(k)^d - a}{(k)^d - b} \frac{1}{b!} m^{\wt(b)}.
    \]
\end{lemma}
\begin{proof}
    It follows from $p_{b,c}^a(m) = \prod_{i=1}^d p_{b_i c_i}^{a_i}(m)$ that $p_{b,c}^a(m)>0$ iff each triple $a_i,b_i,c_i$ satisfies the triangle condition $|b_i-c_i|\leq a_i \leq b_i+c_i$. In this case
    \[
        \deg(p_{b,c}^a(m)) = \sum_{i=1}^d \deg(p_{b_i, c_i}^{a_i}(m))\leq \sum_{i=1}^d\min(b_i,c_i) = \wt(\min(b, c)).
    \]
    Equality holds if and only if $a_i \leq \max(b_i, c_i)$ for all $i$.
    
    If $\deg(p_{b,c}^a(m)) = \wt(b) = \wt(c)$ then $\wt(\min(b, c)) = \wt(b) = \wt(c)$, which implies $b = c$, and moreover we have seen that we must have $a \leq b$.
    Multiplying the leading terms of $p^{a_i}_{b_i,c_i}(m)$ given by the previous lemma, we get the claimed formula.
\end{proof}


Let $\X$ be a fusion of $\cJ(m, k)^d$.
Since $\X$ is a coarsening of $\cJ(m, k)^d$ there is a partition $\fS$ of $[0,k]^d$ such that $\X = \{R_\alpha : \alpha \in \fS\}$, where
\[
    (u, v) \in R_\alpha \iff (|u_1 \setminus v_1|, \dots, |u_d \setminus v_d|) \in \alpha.
\]
In this situation we write $\X = \cJ(m, k)^\fS$ \cite[Section~4]{E}.
For $\beta, \gamma \in \fS$ and $a \in [0,k]^d$ define
\[
    p^a_{\beta,\gamma}(m) = \sum_{b\in \beta, c \in \gamma} p^a_{b,c}(m).
\]
For $\X = \cJ(m,k)^\fS$ to be an association scheme, $\fS$ must satisfy two conditions:
\begin{enumerate}
    \item $\{(0)^d\} \in \fS$,
    \item $p^a_{\beta,\gamma}(m) = p^{a'}_{\beta,\gamma}(m)$ for all $\alpha,\beta,\gamma \in \fS$ and $a,a' \in \alpha$.
\end{enumerate}
We may denote the common value of $p^a_{\beta,\gamma}(m) ~ (a\in\alpha)$ by $p^\alpha_{\beta,\gamma}(m)$; these are the structure constants of $\X$.


We call the sets $\alpha \in \fS$ the \emph{basic $\fS$-sets};
their unions are called \emph{$\fS$-sets}.
For nonempty $S \subseteq [0,k]^d$ let $\wt(S) = \max\{\wt(b) \mid b \in S\}$.
For $\alpha \in \fS$ let $\alpha^*$ be the set of $a \in \alpha$ of maximal weight.
Let $D_\alpha = \bigcup_{a \in \alpha^*} [a]$. Note that $\alpha^*\subseteq D_\alpha$.


{For any $\alpha,\beta\subseteq [0,k]^d$ and $a\in [0,k]^d$ the structure constant $p_{\alpha,\beta}^a(m)$ is a real polynomial in $m$ the coefficients of which depend on $\alpha,\beta$ and $a$. For every pair of distinct real polynomials $f,g\in \mathbb{R}[x]$ there exists a real number $c_{f,g}\in\mathbb{R}$ such that~$f(x)\neq g(x)$ holds for all $x > c_{f,g}$. Therefore, there exists a constant $m_0(k,d)$ such that, provided $m \geq m_0(k, d)$, the following condition holds for all $\alpha,\beta\subset [0,k]^d$ and~$a,b\in [0,k]^d$:
\begin{align}\label{poly}
    p_{\alpha,\beta}^a(m) = p_{\alpha,\beta}^b(m)
    &\implies p_{\alpha,\beta}^a(m) = p_{\alpha,\beta}^b(m)\text{ as polynomials in } m \\
    \notag &\implies \lambda(p_{\alpha,\beta}^a(m)) = \lambda(p_{\alpha,\beta}^b(m)).
\end{align}}

{In what follows we assume that $m \geq m_0(k, d)$ and hence $p^a_{\beta,\gamma}(m) = p^{a'}_{\beta,\gamma}(m)$ only if $p^a_{\beta,\gamma}(m)$ and $p^{a'}_{\beta,\gamma}(m)$ are equal as polynomials in $m$.}
